\scanpage{78}
\mspage{13}
% archive foliation: 80
\uline{Question} Ces flèches sont elles des iso
--- et si oui, comment expliciter un
iso en sens inverse ? \deleted{\unsure{Si}}

Si la détermination conjecturale des
syst. admissibles de disques sur une
surface holomorphe \inserted{compacte} ponctuée (de type
hyperbolique) était connue, on trouve-
rait aisément \deleted{\unsure{des flèches}} un sens
inverse; Au niveau plus naïf des
foncteurs (entre les catégories-solutions
de ces multiplicités) on \unsure{donnerait|fournirait}
ceci : à l'aide d'un \deleted{\unsure{u}}
\[
u_i\in T_{x_i'}^*\wedge_{\mathbb C^*}T_{\unsure{a|x}_i''}^*
\]
\margin{ou ce qui revient au $=$, un
$v_i\,(=u_i^{-1})\in T_{x_i'}^{\vee *}\wedge T_{x_i''}^{\vee *}$}
\deleted{choisissant}\inserted{i.e. un} \deleted{iso} accouplement \unsure{dualisant}
\diagram{\begin{tikzcd}[ampersand replacement=\&]
\deleted{T_{a_i'}^\vee} \arrow[r,draw=none,"{\deleted{\xrightarrow{\sim}}}" description] \& \deleted{T_{a_i''}} \& T_{x_i'}\otimes T_{x_i''} \arrow[r] \& \mathbb C
\end{tikzcd}}
choisissant un \deleted{\unsure{entier}} disque conforme
$D_i'$ en $T_{a_i'}$, \deleted{\unsure{on en déd.}} d'où disque
conforme plein $D_i''$ dans $T_{a_i''}$, on
trouverait \deleted{des} pour le syst. $(D_i')\ (D_i'')$
des \unsure{réalisations} de ces disques comme
